\chapter{Planteamiento de la comparativa: el conjunto de datos}
\label{cap:dataset}

% BLOQUE «Planteamiento de la comparativa» (capitulo 4 de la estructura del tipo 3),
% junto con los capitulos 7 y 8. La norma pide de este bloque: identificar el problema
% concreto, las soluciones alternativas a evaluar, LOS CRITERIOS DE EXITO y LAS MEDIDAS
% QUE SE VAN A TOMAR. De ahi que la variable objetivo (cap. 8) viva aqui y no mas tarde:
% es la definicion de la medida.

% GUION. Fuentes: doc/dataset_info.md, doc/memoria/materiales.md seccion 1,
% TFM-Quantum/ESTADO_ACTUAL.md secciones 1 y 2.
% Figuras: TFM-Quantum/figures/00_anatomia/ (47 disponibles).

\section{Por qué un gemelo digital}
\label{sec:gemelo}
% Fuente: doc/00_Justificaciones.md seccion 1
Tres motivos independientes, y el tercero es el decisivo:
\begin{itemize}
    \item \textbf{cuota}: el plan abierto de IBM da $10$ minutos de QPU al mes;
    \item \textbf{latencia}: las colas hacen inviable un barrido de $20.000$ circuitos;
    \item \textbf{la etiqueta no existe en hardware}: el modelo aprende la diferencia entre el
          resultado ideal y el ruidoso, y un ordenador cuántico \textbf{solo devuelve el
          ruidoso}. Si el valor exacto se pudiera medir, no haría falta mitigar.
\end{itemize}

\section{El procesador simulado}
\label{sec:chip}
\texttt{ibm\_kingston}, familia Heron r2, $156$ qubits, topología \textit{heavy-hex}, puerta de
dos qubits nativa \texttt{cz}. Un qubit está muerto y se excluye del muestreo.
% Figura sugerida: figures/00_anatomia/chip/ (mapa del chip)

\section{Calibración real, no simulada}
\label{sec:calibracion}
\textbf{$42$ días completos} descargados con \texttt{BackendProperties}: T1, T2, error y duración
de \textbf{cada puerta física}, y error de lectura asimétrico, para cada día.

\subsection{Por qué se abandonó la deriva sintética}
El diseño original degradaba T1/T2 de forma progresiva. Hirasaki et al. demuestran que la
deriva real es \textbf{escalonada}: una degradación lineal habría enseñado un patrón inexistente.

\subsection{Lo que solo aparece en datos reales}
% Fuente: doc/memoria/materiales.md seccion 2
Puertas averiadas que \textbf{IBM no repara en $42$ días}, y qubits cuya calibración viola la
cota física $T_2 \leq 2T_1$. Decisión tomada: \textbf{mantenerlas}, porque un
\texttt{gate\_error} de $1{,}0$ es información válida sobre el estado del chip.

\section{Composición y los tres ejes OOD}
\label{sec:ejes}
% Esta seccion es el nucleo del diseno experimental.
\begin{itemize}
    \item \textbf{Tipo}: entrenamiento $100$\,\% circuitos aleatorios; test con seis tipos nunca
          vistos (HEA, TFIM, QAOA, QFT, GHZ, Bernstein-Vazirani).
    \item \textbf{Tamaño}: entrenamiento hasta $11$ qubits; test de $12$ a $15$.
    \item \textbf{Tiempo}: días de calibración \textbf{disjuntos} entre grupos.
\end{itemize}
\textbf{Los tres son independientes}, así que cada fallo se puede atribuir a su causa.

\section{La etiqueta: la batería de observables}
\label{sec:observables}
% Fuente: doc/02_Los_Observables.md (seccion final)
%
% 🔴 CORREGIDO (sep-2026). HAY QUE DISTINGUIR DOS COSAS, y antes se confundian:
%      el DATASET guarda OCHO observables — no se ha regenerado nada;
%      el MODELO entrena sobre CINCO — se retiraron las tres dispersiones (ago-2026).
%    Los cinco vigentes son «las tres bases de Pauli y dos correladores».
%    ⚠️ Explicar POR QUE se retiraron: no fue por redundancia entre ellas, sino porque el
%       ruido es afin (ruidoso = f·exacto + b) y en ese bloque el desplazamiento b vale
%       el 135-149 % de la degradacion que se quiere predecir.
Las tres bases de Pauli, la paridad global y el correlador de vecinos —más las tres
dispersiones, que están en el conjunto de datos pero no son objetivo del modelo.

\subsection{Por qué una batería y no un solo observable}
Para varios tipos de circuito hay observables que valen \textbf{cero por construcción física}
(un GHZ tiene magnetización media nula siempre). Un objetivo escalar único sería una elección
arbitraria que dejaría ciegos a esos tipos.

\subsection{Por qué no cuesta nada}
El simulador calcula el estado \textbf{una vez} y proyecta cada observable sobre él. Medido:
$5$ observables y $100$ tardan lo mismo.

\subsection{Un matiz que hay que declarar}
% 🔴 CIFRA NO CITABLE. Hay dos mediciones en conflicto para la correlacion del bloque de
%    dispersiones: +0,85/+0,77/+0,80 y +0,73/+0,71/+0,68. NO citar ninguna hasta
%    rehacerlo con el protocolo definitivo. Ver doc/memoria/materiales.md §5.
%    Deja de ser bloqueante porque esos tres observables ya no son objetivo.
Las tres dispersiones \textbf{correlacionan fuertemente entre sí}. La justificación de la
batería es de \textbf{simetría de diseño}, no de independencia estadística.

\section{Representación del circuito}
\label{sec:representacion}
El DAG con $25$ dimensiones por nodo: tipo de puerta, ángulo, posición, grado, y
\textbf{telemetría de los qubits físicos concretos} sobre los que actúa ese día.

\textbf{Decisión deliberada:} no se codifica la identidad del qubit, solo su telemetría. Es lo
que permite generalizar a cualquier número de qubits.
\section{Reproducibilidad y sus límites}
\label{sec:reproducibilidad}
% Remite al capitulo 6, no lo desarrolla aqui.
Las semillas dependen solo del índice de la muestra, así que el conjunto es ampliable sin
regenerar. \textbf{Pero la reproducibilidad es del plan, no del bit} —ver
capítulo~\ref{cap:rigor}.
